% !TeX root = ../../MainBook.tex
\section{实变函数概要}\label{sec:Real_Analysis}

读者可能已经注意到，并不是所有的函数都是可积的，例如$\e^x$的能量与功率都是无限的，因此它既不属于能量信号，也不属于功率信号；另一种情况下，函数本身的性质比较差，以至于在Riemann积分的框架下说不清它的积分（或能量）是什么：
\begin{example}[Dirichlet函数]
    \begin{equation}
        \chi_{\mathbb{Q}}(x)=
        \begin{dcases}
            1, & x \in \mathbb{Q},    \\
            0, & x \in \mathbb{Q}^c.
        \end{dcases}
    \end{equation}
    （其中$\mathbb{Q}$表示有理数集，上标$c$表示取补集）
\end{example}
它在Riemann积分的意义下是不可积的，因为它在所有点不连续，但是这样的函数在\textbf{Lebesgue积分}的意义下可积。尽管工程上我们很少考虑这样的函数，但由Lebesgue积分描述的函数空间（见\refsub{Lp_space}）具有很好的性质，这是Riemann积分做不到的。

对于非反常积分，\textbf{Lebesgue积分是Riemann积分的推广}，Riemann可积的函数一定是Lebesgue可积的，并且在相同的区间上积分值相等。

本节就来介绍一些基础的实变函数理论，以便后续学习中使用。由于篇幅原因，我们不得不略去大多数证明（证明可以参考\cite{real_analysis}），不过读者不妨先接受其中的一些重要结论，尤其是有黑体字阐释的定理。

\subsection{*Lebesgue测度与可测函数}\label{sub:lebesgue}

首先，我们给出\textbf{外测度}和\textbf{可测集}的概念。
\begin{definition}[Lebesgue外测度]
    设集合$A\subset\mathbb{R}$，其Lebesgue外测度定义为
    \begin{equation}
        m^*(A)=\inf\braces{\sum_{k=1}^{\infty}l(I_k)\,\middle|\,A\subset\bigcup_{k=1}^{\infty}I_k},
    \end{equation}
    其中$\inf$取自所有覆盖$A$的可数开区间族$\{I_k\}_{k=1}^{\infty}$，$l(I_k)$表示区间$I_k$的长度。
\end{definition}
这个定义看起来不是很直观，但它很好地\textbf{描述了“集合的长度”这一概念}，例如，对于任何一个区间$I$来说，总有$l(I)=m^*(I)$。此外，外测度的以下基本性质符合我们的直观：
\begin{proposition}[外测度的性质]
    \begin{itemize}
        \item 单调性：设$A\subset B$，则$m^*(A)\leq m^*(B)$；
        \item 平移不变性：实数集$A$的平移记为为$A+y=\{x\in \mathbb{R} \mid x-y\in A\}$，满足$m^*(A+y)=m^*(A)$；
        \item 次可数可加性：设$\{E_k\}_{k=1}^{\infty}$是可数集族，则
              \begin{equation}
                  m^*\paren{\bigcup_{k=1}^{\infty}E_k}\leq\sum_{k=1}^{\infty}m^*(E_k).
              \end{equation}
    \end{itemize}
\end{proposition}

次可数可加性使用不等号，是为了涵盖$\{E_k\}_{k=1}^{\infty}$相互重叠的情形。然而，即使各集合互不相交，外测度也未必满足可数可加性，即
\begin{equation}
    m^*\paren{\bigsqcup_{k=1}^{\infty}E_k}\neq\sum_{k=1}^{\infty}m^*(E_k),
\end{equation}
其中$\bigsqcup E_k$表示集族$\{E_k\}$是彼此不交的。要使测度满足可数可加性，集合列中的每个集合都须\textbf{可测}。

\begin{definition}[可测集]
    若集合$E$满足
    \begin{equation}
        \forall A,\quad m^*(A)=m^*(A\cap E)+m^*(A\cap E^c),
    \end{equation}
    则称$E$为\textbf{可测集}。上述条件称为\textbf{Carathéodory条件}或\textbf{卡氏条件}。可测集的\textbf{Lebesgue测度}定义为其Lebesgue外测度，即$m(E)=m^*(E)$。
\end{definition}
尽管不是所有集合都能用测度来描述长度，但可测意味着一个集合边界清晰，可以以较高的精度描述它的“长度”。

Lebesgue测度对不交可数集族具有可数可加性：
\begin{equation}
    m\paren{\bigsqcup_{k=1}^{\infty}E_k}=\sum_{k=1}^{\infty}m(E_k).
\end{equation}

常见集合几乎都是可测的，例如区间及其交、并集。然而，Vitali定理\cite{real_analysis}表明，可借助选择公理构造任意正测度集的不可测子集。

%\begin{proposition}
%    记$\mathbb{R}$上的可测集族为$\mathcal{M} $，则\\
%    \ding{172}空集$\varnothing $，全集$\mathbb{R}$均属于$\mathcal{M} $；\\
%    \ding{173}若$E\in \mathcal{M} $，则$E^c\in \mathcal{M} $；\\
%    \ding{174}可测集的可数族之并仍然是可测集，即
%    $$\{E_k\}_{k=1}^{\infty}\subset \mathcal{M}\implies\bigcup_{k=1}^{\infty}E_k\in \mathcal{M} $$
%    因此，$\mathcal{M} $构成$\mathbb{R}$上的$\sigma$-代数（我们不去定义这个概念，仅仅指出\textbf{可测集在常规的集合运算下是封闭的}）。
%\end{proposition}
%通过以上三条性质可以推知：可测集的可数族的交仍然是可测集；两个可测集的差集仍然是可测集。

\begin{proposition}
    外测度为0的集合都是可测的，称之为\textbf{零测集}。
\end{proposition}
\begin{proof}
    设$m^*(E)=0$，则对于任意集合$A$，$A\cap E,\ A\cup E^c\subset E$，根据外测度的单调性，有$m^*(A\cap E),\ m^*(A\cap E^c)\leq m^*(E)=0$，因此$m^*(A\cap E)=m^*(A\cap E^c)=0$。
    因此，
    \begin{equation}
        m^*(A)=m^*(A\cap E)+m^*(A\cap E^c)=0.
    \end{equation}
\end{proof}

\begin{example}
    单点集是零测集。由测度的可数可加性，可数集是零测集，特别地，$\mathbb{Q}$是可数的，因而是零测集。但是，零测集未必是可数的，例如Cantor集是不可数单点集之并\cite{real_analysis}。
\end{example}

\begin{definition}[几乎处处]
    对于可测集$E$，称一个性质在$E$上\textbf{几乎处处}(almost everywhere, a.e.)成立，如果存在一个零测子集$E_0$，使得该性质对所有$x\in E\setminus E_0$成立。
\end{definition}

\begin{example}[几乎处处相等]
    函数$f,g:E\to\mathbb{R}$在$E$上\textbf{几乎处处相等}，如果存在一个零测子集$E_0$，使得对于所有$x\in E\setminus E_0$，有$f(x)=g(x)$，记为$f=g,\ \aee$。几乎处处相等是一种等价关系。此时，我们\textbf{不区分$f$和$g$}，因为这样的两个函数性质类似，并且我们将看到在$f,g$可积时，它们的积分值相等。
\end{example}

\begin{definition}[可测函数]
    定义在可测集上的函数$f:E\to\mathbb{R}$是\textbf{可测}的，如果
    \begin{equation}
        \forall c\in\mathbb{R},\quad E_c=\{x\in E \mid f(x)<c\}\text{可测}.
    \end{equation}
\end{definition}

可测函数集构成线性空间，并且两个可测函数之积可测，因此可测函数集构成$\mathbb{R}$-代数。

\begin{example}
    基本初等函数在其定义域上都是可测函数；连续函数都是可测函数。
\end{example}

% 这些公理化的定义为理论推带来了方便，但对初学者而言往往难以理解。下面给出可测集和可测函数的一些刻画，以帮助读者快速建立对它们的直观。数学家Littlewood曾说\cite{real_analysis}，只要掌握了以下三个观点，就相当于掌握了可测集与可测函数的核心内容。
%
% 首先，\textbf{可测集可以用区间的有限不交并近似}：
%
% \begin{theorem}
%     设$E$可测，$m(E)<\infty$，则$\forall \epsilon>0$，存在开有界区间族$\{I_k\}_{k=1}^{n}$，记$U=\bigsqcup_{k=1}^{n} I_k$，有
%     \[m(U\setminus E)+m(E\setminus U)<\epsilon\]
% \end{theorem}
%
% 其次，\textbf{可测函数接近于连续函数}：
%
% \begin{theorem}[Lusin定理]
%     设$f:E\to\mathbb{R}$可测，则$\forall \epsilon>0$，存在连续函数$g:E\to\mathbb{R}$及闭集$F\subset E$，使得$m(E\setminus F)<\epsilon$，且在$F$上有$f=g$。
% \end{theorem}
%
在Lebesgue测度的框架下，逐点收敛性可以减弱为几乎处处收敛性：

\begin{definition}[几乎处处收敛]
    函数列$\{f_n:E\to\mathbb{R}\}_{n=1}^{\infty}$在$E$上\textbf{几乎处处收敛}于函数$f:E\to\mathbb{R}$，如果存在一个零测度集$E_0$，使得
    \begin{equation}
        \forall x\in E\setminus E_0,\quad \lim_{n\to\infty}f_n(x)=f(x),
    \end{equation}
    记为$f_n\to f,\ \aee$。
\end{definition}

\begin{example}
    设$E=[0,1]$，并定义
    \begin{equation}
        f_n(x)=(-1)^n\chi_{\mathbb{Q}}(x).
    \end{equation}
    对任意无理数$x\in E\setminus\mathbb{Q}$，恒有$f_n(x)=0$，故$f_n(x)\to 0$；而对任意有理数$x\in E\cap\mathbb{Q}$，数列$\{f_n(x)\}$在$-1$与$1$之间振荡，因而不收敛。由于$E\cap\mathbb{Q}$是零测集，$f_n$几乎处处收敛于零函数，但并不在$E$上逐点收敛。%值得注意的是，不收敛点集虽然测度为零，却在$E$中稠密。
\end{example}
%
% 对于可测函数，若不区分几乎处处相等的函数，则几乎处处收敛就相当于逐点收敛。在这样的约定下，函数列\textbf{在零测集上不好的性质几乎不会影响到全局性质}。
%
% Littlewood的第三个观点是，\textbf{几乎处处收敛的可测函数列近似于一致收敛}：
%
% \begin{theorem}[Egoroff定理]
%     设$m(E)<\infty$，可测函数列$\{f_n:E\to\mathbb{R}\}$几乎处处收敛于$f:E\to\mathbb{R}$，则$\forall \epsilon>0$，存在闭集$F\subset E$，满足$m(E\setminus F)<\epsilon$，且在$F$上$f_n$一致收敛于$f$，记为$f_n\rightrightarrows f$。
% \end{theorem}

\subsection{*Lebesgue积分}\label{sub:integral}

我们首先回顾Riemann积分的定义。

\begin{definition}[Riemann积分]
    设$f:[a,b]\to\mathbb{R}$有界，$P=\{x_0,x_1,\dots,x_n\}$是区间$[a,b]$的一个划分，$a=x_0<x_1<\cdots <x_n=b$，对于每个子区间$[x_{k-1},x_k]$，定义
    \begin{equation}
        m_k=\inf_{x\in [x_{k-1},x_k]}f(x),\quad M_k=\sup_{x\in [x_{k-1},x_k]}f(x).
    \end{equation}
    由此分别定义$f$对$P$的上、下\textbf{Darboux和}：
    \begin{equation}
        U(f,P)=\sum_{k=1}^{n}M_k(x_k-x_{k-1}),\quad L(f,P)=\sum_{k=1}^{n}m_k(x_k-x_{k-1}).
    \end{equation}
    如果
    \begin{equation}
        \inf_{P}U(f,P)=\sup_{P}L(f,P),
    \end{equation}
    则称$f$在$[a,b]$上\textbf{Riemann可积}，并将上述公共值定义为$f$在$[a,b]$上的\textbf{Riemann积分}，为了区分于Lebesgue积分，记为$(R)\int_{a}^{b}f(x)\,\dd x$。
\end{definition}

这个定义实质上是先定义了\textbf{阶梯函数}的Riemann积分，然后通过阶梯函数对被积函数进行上、下逼近，从而定义了被积函数的Riemann积分。阶梯函数的定义如下：

\begin{definition}[阶梯函数]
    如果函数$\varphi:[a,b]\to\mathbb{R}$的定义域可以被划分为有限的子区间族$\{[x_{k-1},x_k]\}_{k=1}^{n}$（严格来说需要处理区间端点，但区间端点属于左侧、右侧的集合均没有影响），使得$\varphi$在每个子区间的内部$(x_{k-1},x_k)$为常数函数，则称$\varphi$是定义在$[a,b]$上的\textbf{阶梯函数}。

    显然阶梯函数Riemann可积，并且
    \begin{equation}
        \int_{a}^{b}\varphi(x)\,\dd x=\sum_{k=1}^{n}\varphi(x)(x_k-x_{k-1}).
    \end{equation}
\end{definition}

Lebesgue定理\cite{real_analysis}表明，闭区间上的有界函数$f$Riemann可积的充要条件是$f$在$[a,b]$中的不连续点构成零测集。Dirichlet函数$\chi_{\mathbb{Q}}$在$[0,1]$处处不连续，因此它在Riemann积分的意义下不可积，而\textbf{Lebesgue积分扩展了有界区间上可积函数的范围}，这是Lebesgue积分的第一条重要优势。

Lebesgue积分的定义与Riemann积分类似，区别在于它需要通过\textbf{简单函数}来逼近可测函数。

\begin{definition}[简单函数]
    若存在函数$\varphi:E\to\mathbb{R}$的定义域$E$的一个有限的可测集划分$\{E_k\}_{k=1}^{n}$，使得$\varphi$在每个可测子集$E_k$上为常数函数，则称$\varphi$是定义在$E$上的\textbf{简单函数}(simple function)。

    简单函数可以表示为：
    \begin{equation}
        \varphi=\sum_{k=1}^{n}a_k\chi_{E_k},\quad \text{其中}E_k=\varphi^{-1}(a_k)=\{x\in E \mid \varphi(x)=a_k\},
    \end{equation}
    称之为简单函数的\textbf{典范表示}(canonical representation)。
\end{definition}

显然，简单函数是可测的。对于非负简单函数$\varphi=\sum_{k=1}^{n}a_k\chi_{E_k}$，定义
\begin{equation}
    \int_{E}\varphi\,\dd m=\sum_{k=1}^{n}a_k m(E_k),
\end{equation}
其中$\dd m$表示Lebesgue测度，并约定$0\cdot\infty=0$。积分值允许为$+\infty$。
\begin{remark}*
    对于单调递增的有界变差函数$f$，可以由$f$的增量定义新的测度，并由此引出\textbf{Lebesgue-Stieltjes积分}，此时$\dd m$就应替换为$\dd f$，但我们不在这里展开讨论，感兴趣的读者可以参考\cite{simple_real_analysis}或\cite{chentianquan}，它实际上是分部积分的基础。
\end{remark}

我们先用非负简单函数从下方逼近非负可测函数来定义其Lebesgue积分，再借助函数的正部和负部处理一般情形。

\begin{definition}[非负可测函数的Lebesgue积分]
    设$f:E\to[0,+\infty)$是非负可测函数，则定义$f$在$E$上的\textbf{Lebesgue积分}为
    \begin{equation}
        \int_{E}f\,\dd m=\sup_{0\leq \varphi\leq f}\int_{E}\varphi\,\dd m,
    \end{equation}
    其中$\varphi$取遍定义在$E$上的非负简单函数。该积分取值于$[0,+\infty]$。

    若$\int_{E}f\,\dd m<\infty$，则称$f$在$E$上\textbf{可积}。不可积的非负函数仍有Lebesgue积分，其值为$+\infty$。
\end{definition}

这个定义不要求$f$有界，也不要求$E$的测度有限。若$m(E)<\infty$且$f$有界，则$f$自然可积。

\begin{theorem}[非负可测函数的Lebesgue积分的线性性]
    设$f,g:E\to[0,+\infty)$是非负可测函数，$\forall\alpha,\beta\in[0,+\infty)$，有
    \begin{equation}
        \int_{E}(\alpha f+\beta g)\,\dd m=\alpha\int_{E}f\,\dd m+\beta\int_{E}g\,\dd m,
    \end{equation}
    其中涉及无穷大的运算按$0\cdot\infty=0$处理。
\end{theorem}

非负可测函数的积分还具有单调性和区域可加性：若$0\leq f\leq g$，则
\begin{equation}
    \int_E f\,\dd m\leq\int_E g\,\dd m;
\end{equation}
若$A,B\subset E$是互不相交的可测集，则
\begin{equation}
    \int_{A\sqcup B}f\,\dd m=\int_A f\,\dd m+\int_B f\,\dd m.
\end{equation}

积分不受零测集上的函数值影响，具体表现为下面的命题。

\begin{proposition}
    设$f:E\to[0,+\infty)$是非负可测函数，则\\
    \ding{172} $\int_{E}f\,\dd m=0\iff f=0,\ \aee$；\\
    \ding{173} 若$E_0$是$E$的零测子集，则$\int_{E}f\,\dd m=\int_{E\setminus E_0}f\,\dd m$。
\end{proposition}

为了将积分推广到可能取正值和负值的函数，我们把函数分解为两个非负部分。

\begin{definition}[函数的正部与负部]
    函数$f:E\to\mathbb{R}$的\textbf{正部}与\textbf{负部}分别定义为
    \begin{equation}
        f^+(x)=\max\{f(x),0\},\quad f^-(x)=-\min\{f(x),0\}.
    \end{equation}
    易见$f=f^+-f^-,\ |f|=f^++f^-$。
\end{definition}
后文中将定义函数的翻转为$f^-(x)=f(-x)$，请读者注意它们的区别，正部与负部这一概念及其符号仅在本节中使用。

\begin{definition}[可测函数的Lebesgue积分]
    设$f:E\to\mathbb{R}$是可测函数，如果$f^+,f^-$均可积，则称$f$是可积的，并定义$f$在$E$上的Lebesgue积分为
    \begin{equation}
        \int_{E}f\,\dd m=\int_{E}f^+\,\dd m-\int_{E}f^-\,\dd m.
    \end{equation}
    $f$的可积性等价于$|f|$的可积性，因为
    \begin{equation}
        \int_{E}|f|\,\dd m=\int_{E}f^+\,\dd m+\int_{E}f^-\,\dd m<\infty.
    \end{equation}
\end{definition}

由定义立即得到Lebesgue积分的绝对值不等式
\begin{equation}
    \abs{\int_{E}f\,\dd m}\leq \int_{E}|f|\,\dd m.
\end{equation}

对于Riemann积分，绝对可积性强于可积性，并称可积却不绝对可积的函数是\textbf{条件可积}的，例如，\cref{def:sinc}给出的采样函数就是条件可积的。但是，从前一定义可以看出，对于Lebesgue积分，\textbf{可积性与绝对可积性是等价的}。

\begin{proposition}[Lebesgue积分的基本性质]
    设$f,g:E\to\mathbb{R}$可积，$\alpha,\beta\in\mathbb{R}$，则：\\
    \ding{172} $\alpha f+\beta g$可积，并且
    \begin{equation}
        \int_E(\alpha f+\beta g)\,\dd m=\alpha\int_E f\,\dd m+\beta\int_E g\,\dd m;
    \end{equation}
    \ding{173} 若$f\leq g$，则
    \begin{equation}
        \int_E f\,\dd m\leq\int_E g\,\dd m;
    \end{equation}
    \ding{174} 若$A,B\subset E$是互不相交的可测集，则
    \begin{equation}
        \int_{A\sqcup B}f\,\dd m=\int_A f\,\dd m+\int_B f\,\dd m.
    \end{equation}
\end{proposition}

\begin{theorem}[比较判别法]
    设$f,g:E\to\mathbb{R}$是可测函数，$g\geq|f|$，如果$g$可积，则$f$可积，并且
    \begin{equation}
        \abs{\int_{E}f\,\dd m}\leq \int_{E}|f|\,\dd m\leq \int_{E}g\,\dd m.
    \end{equation}
\end{theorem}

由于单点集是零测集，区间端点是否包含在积分区域中不会影响积分值，因此记号$\int_{a}^{b}f\,\dd m$是良定义的：
\begin{equation}
    \int_{a}^{b}f\,\dd m=\int_{[a,b]}f\,\dd m=\int_{(a,b)}f\,\dd m=\int_{[a,b)}f\,\dd m=\int_{(a,b]}f\,\dd m.
\end{equation}
在不致混淆时，也沿用$\int_{a}^{b}f(x)\,\dd x$的记号表示Lebesgue积分。

\begin{theorem}
    设$f:[a,b]\to\mathbb{R}$是Riemann可积函数，则$f$也是Lebesgue可积的，并且
    \begin{equation}
        (R)\int_{a}^{b}f(x)\,\dd x=\int_{a}^{b}f\,\dd m.
    \end{equation}
\end{theorem}

可见，在有界区间上，\textbf{Lebesgue积分是Riemann积分的推广}。

下面依次介绍三个用于交换极限与积分顺序的重要结果。前两个定理针对非负函数，控制收敛定理则适用于一般可积函数列。

\begin{theorem}[Levi单调收敛定理]\label{thm:MCT}
    设$\{f_n:E\to[0,+\infty)\}$是递增的非负可测函数列，且$f_n\to f,\ \aee$，则
    \begin{equation}
        \int_{E}f\,\dd m=\lim_{n\to\infty}\int_{E}f_n\,\dd m.
    \end{equation}
\end{theorem}
后文中将它简称为“单调收敛定理”。它表明，对于非负递增的可测函数列，总能\textbf{交换极限与积分}的顺序。特别地，若$\sum_{n=1}^{\infty}f_n$是非负可测函数项级数，其部分和$s_n=\sum_{i=1}^{n}f_i$几乎处处收敛于$s$，则
\begin{equation}
    \int_{E}s\,\dd m=\lim_{n\to\infty}\int_{E}s_n\,\dd m=\sum_{n=1}^{\infty}\int_{E}f_n\,\dd m.
\end{equation}
因此，非负可测函数项级数可以\textbf{逐项积分}。

\begin{corollary}[离散单调收敛定理]\label{cor:discrete_MCT}
    设非负序列$\{a_k[n]\}_{k=1}^{\infty}$满足对每个$n\in\mathbb{Z}$都有$0\leq a_k[n]\leq a_{k+1}[n]$且$\lim_{k\to\infty}a_k[n]=a[n]$，则
    \begin{equation}
        \sum_{n=-\infty}^{\infty}a[n]=\lim_{k\to\infty}\sum_{n=-\infty}^{\infty}a_k[n].
    \end{equation}
    特别地，对单边序列（指标$n\in\mathbb{N}$或$n\in\mathbb{N}^*$），令其余指标处取值为$0$即得相应结论。
\end{corollary}
\begin{proof}
    在$\mathbb{R}$上构造分段常值非负可测函数列
    \begin{equation}
        f_k(t)=\sum_{n=-\infty}^{\infty}a_k[n]\chi_{[n,n+1)}(t),\quad f(t)=\sum_{n=-\infty}^{\infty}a[n]\chi_{[n,n+1)}(t).
    \end{equation}
    显然$\{f_k\}$是非负递增可测函数列，且逐点收敛于$f$。由非负可测函数项级数的逐项积分性质，有
    \begin{equation}
        \int_{\mathbb{R}}f_k\,\dd m=\sum_{n=-\infty}^{\infty}a_k[n],\quad \int_{\mathbb{R}}f\,\dd m=\sum_{n=-\infty}^{\infty}a[n].
    \end{equation}
    对$\{f_k\}$应用单调收敛定理（\cref{thm:MCT}），即得
    \begin{equation}
        \sum_{n=-\infty}^{\infty}a[n]=\int_{\mathbb{R}}f\,\dd m=\lim_{k\to\infty}\int_{\mathbb{R}}f_k\,\dd m=\lim_{k\to\infty}\sum_{n=-\infty}^{\infty}a_k[n].
    \end{equation}
\end{proof}

\begin{theorem}[Fatou引理]
    设$\{f_n:E\to[0,+\infty)\}$是非负可测函数列，则
    \begin{equation}
        \int_{E}\liminf_{n\to\infty}f_n\,\dd m\leq \liminf_{n\to\infty}\int_{E}f_n\,\dd m.
    \end{equation}
    特别地，若$f_n\to f,\ \aee$，则
    \begin{equation}
        \int_{E}f\,\dd m\leq \liminf_{n\to\infty}\int_{E}f_n\,\dd m.
    \end{equation}
\end{theorem}

Fatou引理给出了积分与下极限之间的基本不等式，也是证明其他收敛定理的重要工具。

下面的定理是本书中\textbf{用于极限与积分换序的主要定理}：
\begin{theorem}[Lebesgue控制收敛定理]
    \label{thm:DCT}
    设$\{f_n:E\to\mathbb{R}\}$是可测函数列，$f_n\to f,\ \aee$，如果存在一个可积函数$g:E\to[0,+\infty)$，使得$\forall n\in\mathbb{N}$，$\forall x\in E$，有$|f_n(x)|\leq g(x)$，则$f$可积，并且
    \begin{equation}
        \int_{E}f\,\dd m=\lim_{n\to\infty}\int_{E}f_n\,\dd m.
    \end{equation}
\end{theorem}

后文中将它简称为“控制收敛定理”。\textbf{此定理容易推广到实变量的复值函数列}，只要分别考虑实部和虚部即可。

\begin{corollary}[离散控制收敛定理]\label{cor:discrete_DCT}
    设复值序列$\{a_k[n]\}_{k=1}^{\infty}$对每个$n\in\mathbb{Z}$都满足$\lim_{k\to\infty}a_k[n]=a[n]$。如果存在非负序列$g[n]$满足$\sum_{n=-\infty}^{\infty}g[n]<\infty$，使得$\forall k\in\mathbb{N}^*,\ \forall n\in\mathbb{Z}$都有$|a_k[n]|\leq g[n]$，则级数$\sum_{n=-\infty}^{\infty}a[n]$绝对收敛，并且
    \begin{equation}
        \sum_{n=-\infty}^{\infty}a[n]=\lim_{k\to\infty}\sum_{n=-\infty}^{\infty}a_k[n].
    \end{equation}
    特别地，对单边序列（指标$n\in\mathbb{N}$或$n\in\mathbb{N}^*$），令其余指标处取值为$0$即得相应结论。
\end{corollary}
\begin{proof}
    在$\mathbb{R}$上构造分段常值可测函数列与控制函数：
    \begin{align}
        f_k(t)&=\sum_{n=-\infty}^{\infty}a_k[n]\chi_{[n,n+1)}(t),\quad f(t)=\sum_{n=-\infty}^{\infty}a[n]\chi_{[n,n+1)}(t), \nonumber\\
        G(t)&=\sum_{n=-\infty}^{\infty}g[n]\chi_{[n,n+1)}(t).
    \end{align}
    显然$f_k(t)\to f(t)$逐点成立，且$|f_k(t)|\leq G(t),\ \forall k\in\mathbb{N}^*,\ \forall t\in\mathbb{R}$。由非负可测函数项级数的逐项积分性质，有
    \begin{equation}
        \int_{\mathbb{R}}G\,\dd m=\sum_{n=-\infty}^{\infty}g[n]<\infty,
    \end{equation}
    故$G$在$\mathbb{R}$上可积。又因部分和$f_k^{(N)}(t)=\sum_{n=-N}^{N}a_k[n]\chi_{[n,n+1)}(t)$满足$f_k^{(N)}\to f_k$且$|f_k^{(N)}|\leq G$，由控制收敛定理（\cref{thm:DCT}）可得
    \begin{equation}
        \int_{\mathbb{R}}f_k\,\dd m=\lim_{N\to\infty}\sum_{n=-N}^{N}a_k[n]=\sum_{n=-\infty}^{\infty}a_k[n],
    \end{equation}
    同理$\int_{\mathbb{R}}f\,\dd m=\sum_{n=-\infty}^{\infty}a[n]$（且$|a[n]|\leq g[n]$保证其绝对收敛）。最后对$\{f_k\}$应用控制收敛定理（\cref{thm:DCT}），即得
    \begin{equation}
        \sum_{n=-\infty}^{\infty}a[n]=\int_{\mathbb{R}}f\,\dd m=\lim_{k\to\infty}\int_{\mathbb{R}}f_k\,\dd m=\lim_{k\to\infty}\sum_{n=-\infty}^{\infty}a_k[n].
    \end{equation}
\end{proof}

作为使用控制收敛定理的例子，也作为Lebesgue积分的优秀性质，我们给出以下定理：
\begin{theorem}[Lebesgue积分的可数可加性]
    设$f:E\to\mathbb{R}$是可测函数，$\{E_k\}_{k=1}^{\infty}$是$E$的可测划分，则
    \begin{equation}
        \int_{E}f\,\dd m=\sum_{k=1}^{\infty}\int_{E_k}f\,\dd m.
    \end{equation}
\end{theorem}
\begin{proof}
    记$\chi_n$为$\sqcup_{k=1}^n E_k$的特征函数，则$f\cdot\chi_n\to f$，且$|f\cdot\chi_n|\leq |f|$，由控制收敛定理，
    \begin{equation}
        \int_{E}f\,\dd m=\lim_{n\to\infty}\int_{E}f\cdot\chi_n\,\dd m=\lim_{n\to\infty}\sum_{k=1}^{\infty}\int_{E_k}f\cdot\chi_n\,\dd m=\sum_{k=1}^{\infty}\int_{E_k}f\,\dd m.
    \end{equation}
\end{proof}
\begin{corollary}
    设$f:E\to\mathbb{R}$可积，则\\
    \ding{172} 若$\{A_n\}_{n=1}^{\infty}$是$E$的上升的可测集子族，即$A_n\subset A_{n+1}\subset E,\ \forall n$，则
    \begin{equation}
        \int_{\bigcup_{n=1}^{\infty}A_n}f\,\dd m=\lim_{n\to\infty}\int_{A_n}f\,\dd m;
    \end{equation}
    \ding{173} 若$\{B_n\}_{n=1}^{\infty}$是$E$的下降的可测集子族，即$B_n\supset B_{n+1}\supset E,\ \forall n$，则
    \begin{equation}
        \int_{\bigcap_{n=1}^{\infty}B_n}f\,\dd m=\lim_{n\to\infty}\int_{B_n}f\,\dd m.
    \end{equation}
\end{corollary}
\begin{proof}
    通过一个常见的手法，我们将\textbf{具有单调性的集合列转化为不交的集合列}。令$E_1=A_1,\ E_n=A_n\setminus A_{n-1},\ n\geq 2$，则$\{E_n\}_{n=1}^{\infty}$是不交的，并且$\bigcup_{n=1}^{\infty}A_n=\bigcup_{n=1}^{\infty}E_n$，由积分的可数可加性，
    \begin{equation}
        \int_{\bigcup_{n=1}^{\infty}A_n}f\,\dd m=\sum_{n=1}^{\infty}\int_{E_n}f\,\dd m=\lim_{n\to\infty}\sum_{i=1}^{n}\int_{E_i}f\,\dd m=\lim_{n\to\infty}\int_{\bigcup_{i=1}^{n}A_i}f\,\dd m=\lim_{n\to\infty}\int_{A_n}f\,\dd m.
    \end{equation}
    因此\ding{172}得证，考虑$B_n=E\setminus A_n$并用De Morgan律：
    \begin{equation}
        \bigcup_{n=1}^{\infty}B_n=\bigcup_{n=1}^{\infty}(E\setminus A_n)=E\setminus \bigcap_{n=1}^{\infty}A_n=E,
    \end{equation}
    即可证明\ding{173}。
\end{proof}

特别地，考虑上升的集合列$\{[-n,n]\}_{n=1}^{\infty}$及$\mathbb{R}$上的可积函数$f$，有
\begin{equation}
    \int_{\mathbb{R}}f\,\dd m=\lim_{n\to\infty}\int_{[-n,n]}f\,\dd m.
\end{equation}
这个性质也是Riemann积分不具备的。Riemann积分的无穷限反常积分需要考虑积分上下限是否有关联，例如主值意义下的积分有时不等于普通意义下的积分，而Lebesgue积分则不存在这样的情况。

\begin{example}
    考虑函数$f(x)=x$，它在$\mathbb{R}$上不可积，但在区间$[-n,n]$上可积，并且
    \begin{equation}
        \int_{[-n,n]}f\,\dd m=0.
    \end{equation}
    我们定义$f$在主值意义下的积分为
    \begin{equation}
        \pv\int_{\mathbb{R}}f\,\dd m=\lim_{n\to\infty}\int_{[-n,n]}f\,\dd m=0.
    \end{equation}
    在主值意义下，$f$是可积的。但在Lebesgue积分的意义下，$f$不可积，因为
    \begin{equation}
        \int_{\mathbb{R}}|f|\,\dd m=\int_{\mathbb{R}}|x|\,\dd m=+\infty.
    \end{equation}
\end{example}

\textbf{Riemann积分的计算技巧在多数情况下同样适用于Lebesgue积分}：
\begin{theorem}[Lebesgue积分的变量代换]
    设$f:[a,b]\to\mathbb{R}$可积，递增函数$g:[c,d]\to \mathbb{R}$可导，并且$g(c)=a,\ g(d)=b$，则
    \begin{equation}
        \int_{a}^{b}f\,\dd m=\int_{c}^{d}f\circ g\cdot g'\,\dd m.
    \end{equation}
\end{theorem}

\begin{theorem}[Lebesgue积分的分部积分法]
    设$f,g:[a,b]\to\mathbb{R}$可导，则
    \begin{equation}
        \int_{a}^{b}f g'\,\dd m=\evalat{fg}{a}{b}-\int_{a}^{b}f' g\,\dd m.
    \end{equation}
\end{theorem}

\subsection{\texorpdfstring{$L^p$空间与$l^p$空间}{Lp空间与lp空间}}\label{sub:Lp_space}

本小节主要介绍两类常用的线性赋范空间：$L^p$空间和$l^p$空间（$1\leq p\leq \infty$），前者是处理连续信号的，而后者是处理序列的。

\begin{definition}[$L^p$空间]\label{def:Lp}
    设$E$是可测集，$L^p(E)\ (0<p<\infty)$空间表示在$E$上$p$次Lebesgue可积的函数组成的函数空间，即
    \begin{equation}
        L^p(E):=\braces{f:E\rightarrow \mathbb{C} \,\middle|\, \int_{E}|f(t)|^p\,\dd m<\infty}.
    \end{equation}
\end{definition}

\begin{definition}[$L^{\infty}$空间]\label{def:L_infty}
    设$E$是可测集，$L^{\infty}(E)$空间表示在$E$上\textbf{本性有界}(essentially bounded)的函数组成的函数空间，即
    \begin{equation}
        L^{\infty}(E):=\braces{f:E\rightarrow \mathbb{C} \,\middle|\, \exists M<\infty,\ |f|\leq M,\ \aee}.
    \end{equation}
\end{definition}

\begin{definition}[$l^p$空间]\label{def:lp}
    设$E$是可数集，$l^p(E)\ (0<p<\infty)$空间表示在可数集$E$上$p$次可和的函数组成的函数空间，即
    \begin{equation}
        l^p(E):=\braces{f:E\rightarrow \mathbb{C} \,\middle|\, \sum_{n\in E}|f(n)|^p<\infty}.
    \end{equation}
\end{definition}

一般而言，在考虑$l^p$空间时，可数集$E$取为$\mathbb{N}$，$\mathbb{N}^*$（数列）或$\mathbb{Z}$（序列）。本小节的后续内容中，默认取$E=\mathbb{N}$。

\begin{definition}[$l^{\infty}$空间]\label{def:l_infty}
    $l^{\infty}$表示有界数列空间，即
    \begin{equation}
        l^{\infty}:=\braces{\{x_n\}_{n=0}^{\infty} \,\middle|\, \exists M<\infty,\ |x_n|\leq M,\ \forall n\in \mathbb{N}}.
    \end{equation}
\end{definition}

\begin{lemma}
    $\forall a,b\in\mathbb{C}$，当$p\geq 1$时，有
    \begin{equation}
        |a+b|^p\leq 2^{p-1}(|a|^p+|b|^p)\leq 2^p\max\{|a|^p,|b|^p\},
    \end{equation}
    当$0<p<1$时，有
    \begin{equation}
        |a+b|^p\leq |a|^p+|b|^p.
    \end{equation}
\end{lemma}
\begin{proof}
    由三角不等式，$|a+b|\leq |a|+|b|$。当$p\geq 1$时，函数$x\mapsto x^p$在$[0,\infty)$上是凸函数，因此
    \begin{equation}
        \paren{\frac{|a|+|b|}{2}}^p\leq \frac{|a|^p+|b|^p}{2}.
    \end{equation}
    从而
    \begin{equation}
        |a+b|^p\leq (|a|+|b|)^p\leq 2^{p-1}(|a|^p+|b|^p)\leq 2^p\max\{|a|^p,|b|^p\}.
    \end{equation}

    当$0<p<1$时，若$|a|+|b|=0$，结论显然成立；否则令
    \begin{equation}
        \lambda=\frac{|a|}{|a|+|b|}\in[0,1].
    \end{equation}
    由于$0<p<1$，有$\lambda^p\geq\lambda$及$(1-\lambda)^p\geq 1-\lambda$，故
    \begin{align}
        |a|^p+|b|^p
        &=(|a|+|b|)^p\bigl(\lambda^p+(1-\lambda)^p\bigr)\nonumber\\
        &\geq (|a|+|b|)^p
        \geq |a+b|^p.
    \end{align}
\end{proof}

由此可以推出以上几种空间都是线性空间。此外，在$0<p<1$时，$L^p(E)$和$l^p$空间是度量空间，它们的度量分别为
\begin{equation}
    d(f,g)=\int_{E}|f-g|^p\,\dd m,\quad d(\{x_n\},\{y_n\})=\sum_{n=0}^{\infty}|x_n-y_n|^p.
\end{equation}

不过，$0<p<1$的情况不是我们的主要目标，因为它们不构成范数空间。下面，我们给出以上几种空间在$1\leq p \leq \infty$时的范数定义；由Minkowski不等式可知，这些范数满足三角不等式，因而相应的空间是线性赋范空间。

\begin{definition}[$L^p$空间的范数]
    设$E$是可测集，如果$1\leq p<\infty$，定义
    \begin{equation}
        \|f\|_p=\paren{\int_{E}|f(t)|^p\,\dd m}^{1/p},
    \end{equation}
    如果$p=\infty$，定义
    \begin{equation}
        \|f\|_{\infty}=\inf\{M<\infty \mid |f|\leq M,\ \aee\},
    \end{equation}
    即$f$的本性上确界。
\end{definition}

在$L^p$空间中，我们\textbf{不区分几乎处处相等的函数}，即如果$f=g,\ \aee$，则$f$与$g$在$L^p$空间中是等价的，因此，以上范数定义中的函数$f$可以理解为等价类中的任意一个代表元。在这个意义下，范数的同一性是成立的。容易验证，非负性和齐次性也是成立的。

\begin{definition}[$l^p$空间的范数]
    设$1\leq p<\infty$，对于$\{x_n\}\in l^p$，定义
    \begin{equation}
        \|\{x_n\}\|_p=\paren{\sum_{n=0}^{\infty}|x_n|^p}^{1/p},
    \end{equation}
    对于$\{x_n\}\in l^{\infty}$，定义
    \begin{equation}
        \|\{x_n\}\|_{\infty}=\sup_{n\in \mathbb{N}}|x_n|.
    \end{equation}
\end{definition}

以下设$p,q\in[1,\infty]$满足
\begin{equation}
    \frac{1}{p}+\frac{1}{q}=1,
\end{equation}
其中约定$1/\infty=0$。这样的$p,q$称为一对共轭指数。

\begin{theorem}[Hölder不等式]\label{thm:holder}
    设$f\in L^p(E),\ g\in L^q(E)$，有$f\cdot g\in L^1(E)$，并且
    \begin{equation}
        \int_{E}|fg|\,\dd m\leq \|f\|_p\cdot \|g\|_q.
    \end{equation}
    当$1<p,q<\infty$且$f,g$不几乎处处为零时，等号成立当且仅当存在$\lambda>0$使得$|f|^p=\lambda |g|^q,\ \aee$。
\end{theorem}

\begin{theorem}[Minkowski不等式]\label{thm:minkowski}
    $\forall f,g\in L^p(E)$，有$f+g\in L^p(E)$，并且
    \begin{equation}
        \|f+g\|_p\leq \|f\|_p+\|g\|_p.
    \end{equation}
    当$1<p<\infty$且$f,g$不几乎处处为零时，等号成立当且仅当存在$\lambda>0$使得$f=\lambda g,\ \aee$。
\end{theorem}

Minkowski不等式表明，$L^p(E)$空间的范数满足三角不等式，因此是线性赋范空间。

利用积分的连续性，只要令$\mathbb{R}$上的函数在区间$[n,n+1)\ (n\in\mathbb{Z})$上取常值，即可得到上述不等式的离散形式，从而证明$l^p$空间也是线性赋范空间。

另外，Hölder不等式还能够推出有限测度集上$L^p$空间的包含关系：
\begin{corollary}
    设$m(E)<\infty$，则$\forall 1\leq p_1<p_2\leq \infty$，有$L^{p_2}(E)\subset L^{p_1}(E)$，并且
    \begin{equation}
        \exists c\geq 0,\ \forall f\in L^{p_2}(E),\quad \|f\|_{p_1}\leq c\cdot \|f\|_{p_2}.
    \end{equation}
\end{corollary}
\begin{proof}
    $p_2=\infty$的情况显然成立，下面我们证明$1\leq p_1<p_2<\infty$的情况。

    令$p=p_2/p_1>1,\ 1/p+1/q=1$，则
    \begin{align}
        \|f\|_{p_1}^{p_1} & =\int_{E}|f|^{p_1}\,\dd m\nonumber                                                                      \\
                          & =\int_{E}|f|^{p_1}\cdot \chi_E\,\dd m\nonumber                                                          \\
                          & \leq \paren{\int_{E}|f|^{p_1\cdot p}\,\dd m}^{1/p}\cdot \paren{\int_{E}\chi_E^q\,\dd m}^{1/q}\nonumber \\
                          & =m(E)^{1/q}\cdot \|f\|_{p_2}^{p_1},
    \end{align}
    即
    \begin{equation}
        \|f\|_{p_1}\leq m(E)^{\frac{1}{p_1}-\frac{1}{p_2}}\cdot \|f\|_{p_2}.
    \end{equation}
\end{proof}

对于无限测度集，$L^p$空间互不包含。例如，\refsub{int_sin}中指出采样函数在$\mathbb{R}$上条件可积而不绝对可积，因此$\sin t/t\notin L^1(\mathbb{R})$，但容易验证，$\sin t/t\in L^2(\mathbb{R})$。

$L^p$空间的完备性证明较为复杂，我们仅仅给出结论：
\begin{claim}[Riesz-Fischer定理]\label{claim:Lp_complete}
    设$E$可测，$1\leq p\leq \infty$，则$L^p(E)$空间是完备的。
\end{claim}
容易由$L^p(\mathbb{R})$的完备性推知$l^p$空间的完备性。至此，我们证明了：\textbf{$L^p$空间和$l^p$空间都是完备的线性赋范空间，即Banach空间}。

另外，在$p=2$时，容易给出$L^2(E)$空间的内积（$l^2$空间同理）：
\begin{definition}[$L^2$空间的内积]\label{def:L2_inner_product}
    设$E$可测，定义$L^2(E)$空间的内积为
    \begin{equation}
        (f,g)=\int_{E}f\cdot \overline{g}\,\dd m.
    \end{equation}
\end{definition}

%为了区分于分布理论的符号，本书中总是用$(\cdot,\cdot)$来表示$L^2(E)$空间的内积。

显然，$L^2(E)$空间的范数可以由这个内积导出。因此，\textbf{$L^2(E)$空间构成完备的内积空间，即Hilbert空间}。

令$f=g$，有
\begin{equation}
    (f,f)=\int_{E}|f|^2\,\dd m=E,
\end{equation}
其中$E$为$f$的能量。因此，如果不考虑$(f,f)=0$的情况，则我们在第一节中定义的能量信号恰好就是$L^2(\mathbb{R})$。工程上要求能量信号的能量非零、功率信号的功率非零是自然的，但却带来了数学上的麻烦，因此，\textbf{本书的后续内容中不特意区分能量信号、功率信号的能量和功率是否为0}。

%下面考虑$L^p(E)$空间的对偶空间。由Hölder不等式可知，任取$L^q(E)$中的函数$g$，$g$都定义了一个$L^p(E)$的泛函：\begin{align*}
%    \Phi(f)=\int_{E}f\cdot g\,dm,\forall f\in L^p(E)
%\end{align*}

%\begin{definition}
%    称泛函$\Phi\in [L^p(E)]^*$由函数$g\in L^q(E)$\textbf{表示}(represent)，如果\begin{align*}
%        \forall f\in L^p(E),\Phi(f)=\int_{E}f\cdot g\,dm
%    \end{align*}
%    记为$\Phi=\mathcal{R} (g)$。
%\end{definition}

%事实上，$[L^p(E)]^*$中的元素总能被$g\in L^q(E)$表示：

%\begin{claim}[$L^p(E)$的对偶的Riesz表示定理]
%    设$E$可测，$1\leq p\leq \infty$，则$\forall \Phi\in[L^p(E)]^*$，存在唯一的$g\in L^q(E)$，使得$\Phi=\mathcal{R} (g)$。
%\end{claim}

%可以认为，$\mathcal{R} $是一个从$L^q(E)$到$[L^p(E)]^*$的映射。根据上述定理，$\mathcal{R} $是同构映射，\textbf{线性赋范空间$L^q(E)$与$[L^p(E)]^*$同构}。

\subsection{\texorpdfstring{*$\mathbb{R}^n$上的测度与积分}{*Rn上的测度与积分}}

$\mathbb{R}^n$上的Lebesgue测度可以视作$\mathbb{R}$上的Lebesgue测度的\textbf{乘积测度}。乘积测度空间的完整构造通常要借助Carathéodory-Hahn定理，因此本小节仅引入必要的概念和结论。

\begin{definition}[区间]
    集合$I$是$\mathbb{R}^n$中的\textbf{区间}，如果它是$\mathbb{R}$中的区间族$\{I_k\}_{k=1}^n$的$n$重笛卡尔积：
    \begin{equation}
        I=I_1\times I_2\times\cdots\times I_n.
    \end{equation}
    将区间$I$的\textbf{体积}$\vol(I)$定义为：
    \begin{equation}
        \vol(I)=\prod_{k=1}^n l(I_k),
    \end{equation}
    其中$l(\cdot)$表示$\mathbb{R}$中的区间长度。
\end{definition}

类似于一维的情况，$\mathbb{R}^n$中的外测度定义也由区间的体积出发：

\begin{definition}[$\mathbb{R}^n$上的Lebesgue外测度]
    设集合$A\subset\mathbb{R}^n$，其Lebesgue外测度定义为
    \begin{equation}
        \mu_n^*(A)=\inf\braces{\sum_{k=1}^{\infty}\vol(I_k)\,\middle|\,A\subset\bigcup_{k=1}^{\infty}I_k},
    \end{equation}
    其中$\inf$取自所有覆盖$A$的可数开区间族$\{I_k\}_{k=1}^{\infty}$。
\end{definition}

\begin{definition}[$\mathbb{R}^n$上的Lebesgue可测集]
    若集合$E\subset \mathbb{R}^n$满足卡氏条件：
    \begin{equation}
        \forall A\subset \mathbb{R}^n,\quad \mu_n^*(A)=\mu_n^*(A\cap E)+\mu_n^*(A\cap E^c),
    \end{equation}
    则称$E$为\textbf{$\mu_n$-可测集}，在不引起歧义时也可以直接称之为可测集。可测集的\textbf{Lebesgue测度}定义为其Lebesgue外测度，即$\mu_n(E)=\mu_n^*(E)$。
\end{definition}

$\mu_n(\cdot)$具有与$m(\cdot)$类似的性质，此处不再赘述。特别地，我们仍然约定$0\cdot\infty=0$；作为一种乘积测度，Lebesgue测度是良定义的，即
\begin{equation}
    \mu_{m+n}=\mu_m\times \mu_n=\mu_n\times \mu_m.
\end{equation}

\begin{claim}
    设$L:\mathbb{R}^n\to \mathbb{R}^n$是线性算子，$E\subset\mathbb{R}^n$是$\mu_n$-可测集，则$L(E)$也是$\mu_n$-可测集，并且
    \begin{equation}
        \mu_n(L(E))=|\det L|\mu_n(E).
    \end{equation}
\end{claim}

\begin{example}
    在$\mathbb{R}^n$中，（作为$n$维线性空间的）真子空间$V$的Lebesgue测度为0。这是因为$V$总能在适当的可逆线性变换下变为$\mathbb{R}^k\times\{\mathbf{0}\}$，其中$k$为小于$n$的正整数，根据定义，
    \begin{equation}
        \mu_n(\mathbb{R}^k\times\{\mathbf{0}\})=\vol(\mathbb{R}^k\times\{\mathbf{0}\})=\infty\cdot 0=0.
    \end{equation}
    由此可知$V$的任何子集也都是零测集。
\end{example}

Lebesgue可测函数及其Lebesgue积分的定义全都与\refsub{integral}中一致：我们仍然用简单函数对非负可测函数进行下逼近，然后通过考虑正部与负部推广到一般的可测函数，用于积分与极限换序的单调收敛定理、控制收敛定理也仍然成立，只是函数的定义域更换为$\mathbb{R}^n$。

以下两个定理表明，在多数情况下可以交换重积分的积分次序：
\begin{theorem}[Tonelli定理]\label{thm:tonelli}
    设$X\subset\mathbb{R}^m,\ Y\subset\mathbb{R}^n$，函数$k:X\times Y\to[0,\infty)$是$\mu_{m+n}$-可测的，则
    \begin{enumerate}
        \item 函数$\mathbf{x}\mapsto \int_{Y}k(\mathbf{x},\mathbf{y})\,\dd\mathbf{y}$在$X$上几乎处处有定义且$\mu_m$-可测；
        \item 函数$\mathbf{y}\mapsto \int_{X}k(\mathbf{x},\mathbf{y})\,\dd\mathbf{x}$在$Y$上几乎处处有定义且$\mu_n$-可测；
        \item 积分次序可交换：
              \begin{equation}
                  \int_{X\times Y}k(\mathbf{x},\mathbf{y})\,\dd\mathbf{x}\,\dd\mathbf{y}=\int_X\sbr{\int_Y k(\mathbf{x},\mathbf{y})\,\dd\mathbf{y}}\,\dd\mathbf{x}=\int_Y\sbr{\int_X k(\mathbf{x},\mathbf{y})\,\dd\mathbf{x}}\,\dd\mathbf{y}.
              \end{equation}
    \end{enumerate}
    这里不要求$k$可积，因此以上各式可同时取$\infty$。
\end{theorem}

\begin{theorem}[Fubini定理]\label{thm:fubini}
    设$X\subset\mathbb{R}^m,\ Y\subset\mathbb{R}^n$，函数$k:X\times Y\to\mathbb{R}$可积，则
    \begin{enumerate}
        \item 函数$\mathbf{x}\mapsto \int_{Y}k(\mathbf{x},\mathbf{y})\,\dd\mathbf{y}$在$X$上几乎处处有定义且$\mu_m$-可积；
        \item 函数$\mathbf{y}\mapsto \int_{X}k(\mathbf{x},\mathbf{y})\,\dd\mathbf{x}$在$Y$上几乎处处有定义且$\mu_n$-可积；
        \item 积分次序可交换：
              \begin{equation}
                  \int_{X\times Y}k(\mathbf{x},\mathbf{y})\,\dd\mathbf{x}\,\dd\mathbf{y}=\int_X\sbr{\int_Y k(\mathbf{x},\mathbf{y})\,\dd\mathbf{y}}\,\dd\mathbf{x}=\int_Y\sbr{\int_X k(\mathbf{x},\mathbf{y})\,\dd\mathbf{x}}\,\dd\mathbf{y}.
              \end{equation}
    \end{enumerate}
\end{theorem}
相比Tonelli定理，Fubini定理虽要求$k$可积，却不要求$k$非负。

将序列表示为在各区间$[n,n+1)$上取常值的可测函数，即可由积分换序定理得到相应的级数结论。例如：
\begin{corollary}\label{cor:sequence_fubini_tonelli}
    设$\{a_{ij}\}_{i,j=1}^{\infty}\subset\mathbb{R}$，则当$a_{ij}\ge 0$或$\sum_{j=1}^{\infty}\paren{\sum_{i=1}^{\infty}|a_{ij}|}<\infty$时，可以交换求和次序：
    \begin{equation}
        \sum_{j=1}^{\infty}\paren{\sum_{i=1}^{\infty}a_{ij}}=\sum_{i=1}^{\infty}\paren{\sum_{j=1}^{\infty}a_{ij}}.
    \end{equation}
\end{corollary}

\subsection{正交函数系}\label{sub:orthognal}
\begin{definition}[正交基]
    如果有内积空间$V$的一组基$\{\mathbf{v}_1,\mathbf{v}_2,\dots ,\mathbf{v}_n\}$，满足$\langle \mathbf{v}_i,\mathbf{v}_j\rangle =0,\ i \neq j$，则称这组基是一组\textbf{正交基}(orthogonal bases)。

    如果对$V$中任一向量$\mathbf{w}$，均可以表示为这组基的线性组合：
    \begin{equation}
        \mathbf{w}=\sum_{i = 1}^{n} c_i \mathbf{v}_i,
    \end{equation}
    则这组基是\textbf{完备正交基}(complete orthogonal bases)。

    如果$\{\mathbf{v}_1,\mathbf{v}_2,\dots ,\mathbf{v}_n\}$还满足$\langle \mathbf{v}_i,\mathbf{v}_i\rangle =1,\ i\in \{1,2,\dots ,n\}$，称这组基是\textbf{标准完备正交基}(complete orthonormal bases)。
\end{definition}

\begin{theorem}\label{thm:vector_decomposition}
    设内积空间$V$有完备正交基$\{\mathbf{v}_1,\mathbf{v}_2,\dots ,\mathbf{v}_n\}$，则此空间中任意向量均有分解式
    \begin{equation}
        \mathbf{w}=\sum_{i=1}^{n}\frac{\langle \mathbf{w},\mathbf{v}_j \rangle}{\langle \mathbf{v}_j,\mathbf{v}_j \rangle}\mathbf{v}_i,
    \end{equation}
    特别地，如果$\{\mathbf{v}_1,\mathbf{v}_2,\dots ,\mathbf{v}_n\}$是标准完备正交基，则
    \begin{equation}
        \mathbf{w}=\sum_{i=1}^{n}\langle \mathbf{w},\mathbf{v}_j \rangle\mathbf{v}_i.
    \end{equation}
\end{theorem}
\begin{proof}
    将$\mathbf{w}=\sum_{i = 1}^{n} c_i \mathbf{v}_i$两边同时对$\mathbf{v}_j$做内积，得到：
    \begin{equation}
        \langle \mathbf{w},\mathbf{v}_j \rangle=\braket{\sum_{i = 1}^{n} c_i \mathbf{v}_i,\mathbf{v}_j} =c_j\langle \mathbf{v}_j,\mathbf{v}_j \rangle.
    \end{equation}
    因此
    \begin{equation}
        c_j=\frac{\langle \mathbf{w},\mathbf{v}_j \rangle}{\langle \mathbf{v}_j,\mathbf{v}_j \rangle},\quad \mathbf{w}=\sum_{j = 1}^{n} \frac{\langle \mathbf{w},\mathbf{v}_j \rangle}{\langle \mathbf{v}_j,\mathbf{v}_j \rangle} \mathbf{v}_j.
    \end{equation}
    对于标准正交基，有$\langle \mathbf{v}_j,\mathbf{v}_j \rangle=1$。
\end{proof}

以下定理是勾股定理的高维推广：
\begin{theorem}[毕达哥拉斯定理]\label{thm:gougu}
    设内积空间$V$有完备正交基$\{\mathbf{v}_1,\mathbf{v}_2,\dots ,\mathbf{v}_n\}$，则
    \begin{equation}
        |\mathbf{w}|^2=\sum_{i=1}^{n}|\langle\mathbf{w},\mathbf{v}_i\rangle|^2|\mathbf{v}_i|^2.
    \end{equation}
    特别地，如果$\{\mathbf{v}_1,\mathbf{v}_2,\dots ,\mathbf{v}_n\}$是标准完备正交基，则
    \begin{equation}
        |\mathbf{w}|^2=\sum_{i=1}^{n}|\langle\mathbf{w},\mathbf{v}_i\rangle|^2.
    \end{equation}
\end{theorem}
\begin{proof}
    直接将$\mathbf{w}$与自己做内积并带入\cref{thm:vector_decomposition}，得到
    \begin{align}
        |\mathbf{w}|^2 & =\braket{\sum_{i=1}^{n}\langle \mathbf{w},\mathbf{v}_i\rangle \mathbf{v}_i,\sum_{i=1}^{n}\langle \mathbf{w},\mathbf{v}_i\rangle \mathbf{v}_i}\nonumber \\
                       & =\sum_{i=1}^{n}|\langle\mathbf{w},\mathbf{v}_i\rangle|^2\langle\mathbf{v}_i,\mathbf{v}_i\rangle+\sum_{1\leq i<j\leq n}\langle \mathbf{w},\mathbf{v}_i\rangle\langle \mathbf{w},\mathbf{v}_j\rangle\langle \mathbf{v}_i,\mathbf{v}_j\rangle\nonumber \\
                       & =\sum_{i=1}^{n}|\langle\mathbf{w},\mathbf{v}_i\rangle|^2|\mathbf{v}_i|^2.
    \end{align}
\end{proof}
可见做正交基分解能够极大地简化对线性赋范空间的研究。

我们知道，$L^2([0,T])$空间是完备的内积空间（\cref{claim:Lp_complete}），重述内积的定义\cref{def:L2_inner_product}如下：
\begin{definition}[函数空间上的内积]
    在空间
    \begin{equation}
        L^2([0,T]):=\braces{f:[0,T]\rightarrow \mathbb{C} \,\middle|\, \int_{T}|f(t)|^2\,\dd t<\infty}
    \end{equation}
    上定义内积：
    \begin{equation}
        (f,g):=\int_{T}f(t)\overline{g(t)}\,\dd t,
    \end{equation}
    其中$\int_T$表示在区间$[0,T]$上做积分。对于周期函数，积分区域可以理解为任意长为$T$的区间。
\end{definition}

有了内积，就可以定义正交性，从而有正交函数系的概念：
\begin{definition}[正交函数系]
    在$L^2([0,T])$空间中，如果函数列$\{f_i(t)\}_{i=1}^{\infty}$满足
    \begin{equation}
        (f_i,f_j)=\int_{T}f_i(t)f_j^*(t)\,\dd t=0,\quad i\neq j,
    \end{equation}
    则称$\{f_i(t)\}_{i=1}^{\infty}$为$L^2([0,T])$空间中的一组\textbf{正交函数系}。
\end{definition}

\begin{definition}[完备正交函数系]
    对于一个正交函数系$\{f_i(t)\}_{i=1}^{\infty}$，如果不存在非零的函数$g(t)\notin\{f_i(t)\}_{i=1}^{\infty}$使得$g(t)$与$\{f_i(t)\}_{i=1}^{\infty}$中的所有函数正交，就称$\{f_i(t)\}_{i=1}^{\infty}$为\textbf{完备正交函数系}。
\end{definition}

完备性意味着$L^2([0,T])$空间中的任一函数$f(t)$均可分解为这个函数系的级数$\sum_{i = 1}^{\infty} a_i f_i(t)$，我们沿用\cref{thm:vector_decomposition}计算这些系数：
\begin{equation}
    a_i=\frac{(g,f_i)}{(f_i,f_i)}=\frac{\int_{T}g(t)f_i^*(t)\,\dd t}{\int_{T}|f_i(t)|^2\,\dd t}.
\end{equation}
细心的读者可能已经发现，这里得到的公式用到了有限维线性空间中的结论，但要推广到无限维线性空间并不是显然的。我们将在\refsub{paseval}中给出Parseval恒等式之后一并讨论这个问题。

在下一章中我们将看到，三角函数系和指数函数系都是完备正交函数系，是傅里叶级数理论的基础。在添加适当权函数的前提下，\refsub{chebychev_polynomial}中介绍的Chebyshev多项式以及\refsub{bessel}中介绍的Bessel函数都是完备正交函数系；此外，典型的标准完备正交函数系还有Legendre多项式、小波变换基函数(wavelet basis functions)等。

本书仅讨论三角函数系和指数函数系，它们的本质上都是\refsub{chebychev_polynomial}中介绍的\textbf{三角多项式}的延申。
